\documentclass[11pt, letterpaper]{article}


\usepackage[T1]{fontenc}
\usepackage[utf8]{inputenc}
\usepackage[spanish, es-noshorthands, es-tabla, es-nodecimaldot]{babel}
\usepackage{lmodern}
\usepackage{microtype}
\usepackage[margin=2.5cm]{geometry}

\usepackage{amsmath, amssymb, amsthm, mathtools}

\usepackage{booktabs, tabularx, array}

\usepackage{csquotes}
\usepackage[backend=biber, style=authoryear, maxcitenames=2,
            giveninits=true, uniquename=false, uniquelist=false]{biblatex}
\DeclareDelimFormat{nameyeardelim}{\addcomma\space}

\usepackage{xcolor}
\definecolor{enlace}{RGB}{0, 70, 140}
\usepackage[colorlinks=true, linkcolor=enlace, citecolor=enlace,
            urlcolor=enlace]{hyperref}


\theoremstyle{plain}
\newtheorem{teorema}{Teorema}[section]
\newtheorem{proposicion}[teorema]{Proposición}
\newtheorem{lema}[teorema]{Lema}
\newtheorem{corolario}[teorema]{Corolario}

\theoremstyle{definition}
\newtheorem{definicion}[teorema]{Definición}
\newtheorem{regla}[teorema]{Regla}
\newtheorem{ejemplo}[teorema]{Ejemplo}
\newtheorem{observacion}[teorema]{Observación}

\newcommand{\K}{\mathbb{K}}
\newcommand{\R}{\mathbb{R}}
\DeclareMathOperator{\SSE}{SSE}
\DeclareMathOperator{\AMISE}{AMISE}
\DeclareMathOperator{\IQR}{IQR}

\addbibresource{referencias.bib}

\title{Fundamentos teóricos\\[0.5em]
  \large Análisis Formal de Conceptos para el Análisis de Bases de Datos}
\author{Tomas Jerez \and Juan Huertas \and Martin Hernandez \and Angel Arevalo}
\date{\small Análisis de Bases de Datos\,·\,Ciencias de la Computación\\
  Universidad Nacional de Colombia\\[0.8em]
  \today}

\hypersetup{pdftitle={Fundamentos teóricos}}

\begin{document}

\maketitle

\section*{Sobre este documento}

Este documento reúne las definiciones, resultados y demostraciones en los que se apoya cada fase del proyecto, y crece junto con el trabajo. Cada sección corresponde a una fase del pipeline. Los resultados se numeran por sección (Definición~1.3, Teorema~1.26, \dots) para que las referencias desde el código, los notebooks y el informe no cambien cuando se agregue contenido nuevo.

La explicación de cada módulo de código, con su pseudocódigo, está en la carpeta \path{docs/codigo/} del repositorio. Esos documentos citan los resultados de este por su número.

\paragraph{Convenciones.} Se escribe $\lvert X \rvert$ para el número de elementos de un conjunto~$X$, y $\lfloor y \rfloor$ y $\lceil y \rceil$ para las partes enteras inferior y superior de un real~$y$. En todo el documento $n$ es el número de objetos (filas de la tabla) y $k$ el número de intervalos de una discretización.

\tableofcontents


\section{Del dato tabular al contexto formal}
\label{sec:contexto}

El Análisis Formal de Conceptos trabaja sobre tablas binarias: cada objeto tiene o no tiene cada atributo. Los datos reales traen columnas numéricas (radio, textura, área\dots) y categóricas (diagnóstico). Esta sección formaliza los tres pasos que llevan de una tabla así a un contexto formal: la discretización de las variables numéricas (\S\ref{sec:discretizacion} y \S\ref{sec:numero-intervalos}), el escalado conceptual (\S\ref{sec:escalado} y \S\ref{sec:propiedades-escalado}) y la matriz de incidencia (\S\ref{sec:contextos}). En el código, \path{cuant_type.py} y \path{quantiles.py} discretizan, y \path{scaler.py} escala y produce la matriz; todos están en \path{src/preprocessing/}, y su documentación en \path{docs/codigo/preprocessing/}.


\subsection{Contextos formales y matriz de incidencia}
\label{sec:contextos}

\begin{definicion}[contexto formal]\label{def:contexto}
Un \emph{contexto formal} es una terna $\K = (G, M, I)$ formada por un conjunto finito~$G$ de \emph{objetos}, un conjunto finito~$M$ de \emph{atributos} y una relación $I \subseteq G \times M$. Si $(g, m) \in I$ se escribe $g \mathrel{I} m$: \enquote{el objeto~$g$ tiene el atributo~$m$}.
\end{definicion}

\begin{definicion}[fila y columna]\label{def:fila-columna}
Para $g \in G$ y $m \in M$, $g' = \{ m \in M : g \mathrel{I} m \}$ es la \emph{fila} de~$g$ y $m' = \{ g \in G : g \mathrel{I} m \}$ es la \emph{columna} de~$m$.
\end{definicion}

\begin{definicion}[matriz de incidencia]\label{def:matriz}
Con $G = \{g_1, \dots, g_n\}$ y $M = \{m_1, \dots, m_p\}$ en un orden fijo, la \emph{matriz de incidencia} es $X \in \{0, 1\}^{n \times p}$ con $X_{ij} = 1$ si y solo si $g_i \mathrel{I} m_j$. Se dibuja como una \emph{tabla de cruces}.
\end{definicion}

\begin{observacion}[representación]\label{obs:representacion}
Fijados los órdenes, la matriz determina el contexto. Se puede guardar de forma explícita, como un \texttt{DataFrame} booleano, o implícita, con cada fila~$g'$ como un entero usado como vector de bits; esta es la que usa el código de la Sección~\ref{sec:conceptos}.
\end{observacion}


\subsection{Contextos multivaluados y escalado conceptual}
\label{sec:escalado}

\begin{definicion}[contexto multivaluado]\label{def:multivaluado}
Un \emph{contexto multivaluado} $(G, M, W, I)$ tiene objetos~$G$, atributos multivaluados~$M$, valores~$W$ y una relación $I \subseteq G \times M \times W$ en la que cada par $(g, m)$ tiene a lo sumo un valor, que se escribe $m(g)$; si no tiene ninguno, el valor \emph{falta}. $D_m$ es el conjunto de objetos en los que $m$ está definido y $V_m = \{ m(g) : g \in D_m \}$ el de sus \emph{valores observados}.
\end{definicion}

\begin{observacion}[modelo relacional]\label{obs:relacional}
Un contexto multivaluado es una tabla de una base de datos: los objetos son las tuplas, los atributos son las columnas y un valor faltante es un \texttt{NULL}.
\end{observacion}

\begin{definicion}[escala conceptual]\label{def:escala}
Una \emph{escala} para~$m$ es un contexto formal $\mathbb{S}_m = (S, N_m, I_m)$ cuyos objetos incluyen los valores observados de~$m$; sus atributos~$N_m$ son los atributos binarios con los que se describe cada valor. En el código, $S = V_m$.
\end{definicion}

\begin{definicion}[contexto derivado]\label{def:derivado}
El \emph{contexto derivado} de $(G, M, W, I)$ con escalas~$\mathbb{S}_m$ es el contexto formal $(G, N, J)$ con
\[
  N = \bigcup_{m \in M} \{m\} \times N_m, \qquad g \mathrel{J} (m, a) \iff g \in D_m \text{ y } m(g) \mathrel{I_m} a.
\]
Cada columna~$m$ se reemplaza por un bloque de columnas binarias; si a~$g$ le falta el valor de~$m$, su fila en ese bloque queda vacía.
\end{definicion}

\begin{definicion}[escalas nominal y ordinal]\label{def:escalas}
Para valores $v_1 < \dots < v_r$, la \emph{escala nominal} es $(V, V, =)$, la \emph{ordinal creciente} es $(V, V, \le)$ y la \emph{ordinal decreciente} es $(V, V, \ge)$. En el contexto derivado, sus atributos se escriben $(m{=}w)$, $(m{\le}w)$ y $(m{\ge}w)$, y son las columnas \texttt{m=w}, \texttt{m<=w} y \texttt{m>=w} de la función \texttt{scaling} de \path{scaler.py}.
\end{definicion}

\begin{ejemplo}[tres valores]\label{ej:escalas}
Con $V = \{0, 1, 2\}$, la escala nominal es la matriz identidad, y la ordinal creciente es triangular superior: el valor~0 tiene $(\le 0)$, $(\le 1)$ y $(\le 2)$; el valor~1, $(\le 1)$ y $(\le 2)$; y el valor~2, solo $(\le 2)$.
\end{ejemplo}

\begin{observacion}[escala interordinal]\label{obs:interordinal}
Si se usan a la vez las dos escalas ordinales (la escala \emph{interordinal} de \textcite{ganter1999}), los atributos $(m{\ge}v)$ y $(m{\le}w)$ describen juntos el intervalo $v \le m(g) \le w$. Se obtiene concatenando las salidas de \texttt{scaling} con \texttt{"le"} y con \texttt{"ge"}.
\end{observacion}


\subsection{Propiedades del escalado}
\label{sec:propiedades-escalado}

Sea $m$ un atributo con valores observados $v_1 < \dots < v_r$, escalado sobre~$V_m$ como en el código. Las propiedades siguientes salen directamente de las definiciones, usando que $m(g)$ es único.

\begin{proposicion}[la escala nominal induce una partición]\label{prop:particion}
Los conjuntos $(m{=}v)'$, con $v \in V_m$, son no vacíos, disjuntos dos a dos y su unión es~$D_m$. Sin valores faltantes, cada fila tiene exactamente un~1 en el bloque de~$m$.
\end{proposicion}

\begin{proposicion}[la escala ordinal induce una cadena]\label{prop:cadena}
$(m{\le}v_1)' \subseteq \dots \subseteq (m{\le}v_r)' = D_m$ y $D_m = (m{\ge}v_1)' \supseteq \dots \supseteq (m{\ge}v_r)'$.
\end{proposicion}

\begin{corolario}[atributo universal]\label{cor:universal}
Sin valores faltantes, todos los objetos tienen $(m{\le}v_r)$ y $(m{\ge}v_1)$: sus columnas están llenas. Un atributo así no cambia el retículo y se puede quitar (Proposición~\ref{prop:universal}).
\end{corolario}

\begin{observacion}[regularidades de la escala]\label{obs:fondo}
Por las Proposiciones \ref{prop:particion} y~\ref{prop:cadena}, todo objeto con $(m{\le}v_i)$ tiene $(m{\le}v_j)$ si $i < j$, y ninguno tiene a la vez $(m{=}v)$ y $(m{=}w)$ con $v \ne w$. Son regularidades de la escala, no de los datos; la \S\ref{sec:fondo} las separa de las implicaciones que sí vienen de los datos.
\end{observacion}

\begin{proposicion}[el escalado no pierde información]\label{prop:sin-perdida}
Para $g \in D_m$, el valor $m(g)$ se recupera de la fila~$g'$: es el único~$v$ con $(m{=}v) \in g'$, el menor~$v$ con $(m{\le}v) \in g'$ o el mayor~$v$ con $(m{\ge}v) \in g'$, según la escala.
\end{proposicion}

\begin{corolario}\label{cor:filas}
Dos objetos sin valores faltantes tienen la misma fila en el contexto derivado si y solo si tienen los mismos valores en todas las variables. El contexto derivado tiene $\sum_{m} \lvert V_m \rvert$ atributos.
\end{corolario}

Así, toda la pérdida de información del pipeline ocurre en la discretización. Con las 30~variables de Breast Cancer Wisconsin y $k = 11$ intervalos por variable, el contexto tiene a lo sumo 330~atributos.


\subsection{Discretización}
\label{sec:discretizacion}

Una variable continua toma casi tantos valores distintos como objetos hay, y escalarla directamente daría casi una columna por objeto. La discretización agrupa los valores en pocos intervalos ordenados.

\begin{definicion}[discretización por puntos de corte]\label{def:discretizacion}
Dados cortes $c_0 < c_1 < \dots < c_k$ con $c_0 \le \min x$ y $c_k \ge \max x$, los intervalos son $I_0 = [c_0, c_1]$ e $I_j = (c_j, c_{j+1}]$ para $1 \le j < k$, y la discretización asigna a cada objeto el índice $\delta(g) = j$ del intervalo que contiene a~$x(g)$. Es la convención de pandas.
\end{definicion}

\begin{proposicion}[la discretización conserva el orden]\label{prop:orden}
Si $x(g) \le x(h)$, entonces $\delta(g) \le \delta(h)$.
\end{proposicion}

Por eso se puede aplicar una escala ordinal a los códigos, y cada atributo ordinal se lee en la variable original: $\delta(g) \le j$ si y solo si $x(g) \le c_{j+1}$. Con los cortes del radio de la \S\ref{sec:ejemplo}, por ejemplo, \texttt{radio<=1} significa \enquote{radio${}\le 17.13$}. La discretización es el único paso con pérdida: los objetos cuyos valores caen en un mismo intervalo se vuelven indistinguibles en esa variable. El proyecto ofrece tres formas de elegir los cortes.

\begin{definicion}[igual ancho]\label{def:ancho}
$c_j = a + j(b - a)/k$, con $a = \min x$ y $b = \max x$. Todos los intervalos miden lo mismo, pero un valor atípico estira el rango y puede dejar casi todos los objetos en uno o dos intervalos.
\end{definicion}

\begin{definicion}[igual frecuencia]\label{def:frecuencia}
Los cortes son los cuantiles empíricos $c_j = \hat{Q}(j/k)$, con la interpolación lineal que usan por defecto NumPy y pandas, el tipo~7 de \textcite{hyndman1996}. Con $k = 4$, los cortes interiores son los cuartiles.
\end{definicion}

\begin{proposicion}[intervalos balanceados]\label{prop:balance}
Si los $n$~valores son distintos, cada intervalo de igual frecuencia contiene, en aritmética exacta, $\lfloor n/k \rfloor$ o $\lceil n/k \rceil$ objetos.
\end{proposicion}

\begin{observacion}[empates]\label{obs:empates}
Con valores repetidos, varios cuantiles pueden coincidir; la implementación elimina los cortes repetidos, así que pueden quedar menos de $k$~intervalos. Por redondeo, algún intervalo puede diferir en un objeto del tamaño de la Proposición~\ref{prop:balance}.
\end{observacion}

\begin{definicion}[$k$-medias en una dimensión]\label{def:kmedias}
Se buscan grupos $C_1, \dots, C_k$ y centroides $\mu_1, \dots, \mu_k$ que minimicen $\SSE(C, \mu) = \sum_{i} \sum_{x \in C_i} (x - \mu_i)^2$. El \emph{algoritmo de Lloyd} \parencite{lloyd1982} alterna dos pasos hasta que los centroides no cambian: asignar cada valor al centroide más cercano y reemplazar cada centroide por la media de su grupo.
\end{definicion}

\begin{lema}[la media minimiza la suma de cuadrados]\label{lem:media}
Si $S$ es un multiconjunto finito de reales con media~$\bar{s}$, para todo $c \in \R$ se tiene $\sum_{x \in S} (x - c)^2 = \sum_{x \in S} (x - \bar{s})^2 + \lvert S \rvert \, (\bar{s} - c)^2$.
\end{lema}

\begin{teorema}[convergencia del algoritmo de Lloyd]\label{teo:lloyd}
La SSE no aumenta en ningún paso, disminuye estrictamente en cada iteración en que cambia algún centroide, y los centroides dejan de cambiar después de un número finito de iteraciones.
\end{teorema}

La asignación no aumenta ningún término de la SSE, la actualización tampoco, por el Lema~\ref{lem:media}, y como hay finitas formas de repartir los valores, la SSE no puede bajar indefinidamente.

\begin{observacion}[óptimo local]\label{obs:optimo-local}
El resultado es un óptimo local, que depende de los centroides iniciales; la implementación los pone en percentiles fijos, así que es reproducible. En una dimensión, el óptimo global se puede calcular con programación dinámica \parencite{fisher1958, wang2011}.
\end{observacion}

\begin{proposicion}[regiones de asignación]\label{prop:voronoi}
Si $\mu_1 < \dots < \mu_k$, el centroide más cercano a~$x$ es~$\mu_j$ si y solo si $x$ está entre los puntos medios $(\mu_{j-1} + \mu_j)/2$ y $(\mu_j + \mu_{j+1})/2$.
\end{proposicion}

Por eso \texttt{kmeans\_1d}, al terminar, corta en los puntos medios entre centroides consecutivos, y los códigos conservan el orden (Proposición~\ref{prop:orden}).


\subsection{Número de intervalos}
\label{sec:numero-intervalos}

Con pocos intervalos se pierden distinciones entre objetos (Corolario~\ref{cor:filas}); con muchos, el contexto tiene más atributos y el número de conceptos puede crecer mucho, que es el principal reto del dataset de Spotify. Las reglas siguientes vienen de la teoría de histogramas.

\begin{regla}[\cite{sturges1926}]\label{regla:sturges}
$k = \lceil \log_2 n \rceil + 1$. Sturges idealizó el histograma de una variable normal como una binomial con $k$~intervalos, cuyas frecuencias suman $2^{k-1} = n$. Con $n = 569$ da $k = 11$.
\end{regla}

\begin{teorema}[\cite{scott1979}]\label{teo:amise}
Si la densidad~$f$ es suave y $R(f') = \int f'(x)^2 \, dx$ es finito, el error cuadrático medio integrado de un histograma de ancho~$h$ con $n$~datos es, asintóticamente, $\AMISE(h) = 1/(nh) + h^2 R(f')/12$.
\end{teorema}

\begin{corolario}[ancho óptimo]\label{cor:ancho-optimo}
La AMISE se minimiza en $h^{*} = \bigl( 6 / (n R(f')) \bigr)^{1/3}$, que decrece como $n^{-1/3}$; la regla de Sturges, en cambio, crece solo como $\log_2 n$.
\end{corolario}

\begin{corolario}[regla de Scott]\label{cor:scott}
Para una normal con desviación~$\sigma$, $R(f') = 1/(4\sqrt{\pi}\,\sigma^3)$ y $h^{*} = (24\sqrt{\pi})^{1/3} \sigma\, n^{-1/3} \approx 3.49\, \sigma\, n^{-1/3}$. La implementación estima~$\sigma$ con la desviación estándar muestral.
\end{corolario}

\begin{regla}[\cite{freedman1981}]\label{regla:fd}
$h = 2 \IQR\, n^{-1/3}$, con $\IQR$ el rango intercuartílico. Es la variante robusta de la regla de Scott: el IQR, a diferencia de la desviación estándar, no se ve afectado por los valores atípicos.
\end{regla}

\begin{observacion}[del ancho al número de intervalos]\label{obs:ancho-k}
Con un ancho~$h$, $k = \lceil (\max x - \min x)/h \rceil$; si la dispersión es nula se usa la regla de Sturges, y \texttt{discretize} obliga a que $k \ge 2$. Las reglas se diseñaron para intervalos de igual ancho: usarlas con igual frecuencia o con $k$-medias es una heurística.
\end{observacion}


\subsection{Ejemplo completo}
\label{sec:ejemplo}

Se toman seis objetos con dos variables numéricas y una categórica (valores ilustrativos, inspirados en Breast Cancer Wisconsin). La tabla de la izquierda tiene los valores originales y la de la derecha sus códigos después de discretizar.

\begin{center}
\small
\begin{tabular}{cccc}
  \toprule
  Objeto & radio & textura & diagnóstico \\
  \midrule
  $g_1$ & 12.0 & 15.7 & B \\
  $g_2$ & 13.5 & 21.3 & B \\
  $g_3$ & 20.1 & 17.8 & M \\
  $g_4$ & 11.2 & 25.0 & B \\
  $g_5$ & 17.8 & 19.9 & M \\
  $g_6$ & 14.6 & 14.1 & B \\
  \bottomrule
\end{tabular}
\qquad\qquad
\begin{tabular}{cccc}
  \toprule
  Objeto & radio & textura & diagnóstico \\
  \midrule
  $g_1$ & 0 & 0 & B \\
  $g_2$ & 0 & 1 & B \\
  $g_3$ & 2 & 1 & M \\
  $g_4$ & 0 & 2 & B \\
  $g_5$ & 2 & 1 & M \\
  $g_6$ & 1 & 0 & B \\
  \bottomrule
\end{tabular}
\end{center}

Con intervalos de igual ancho y $k = 3$, los cortes del radio son 11.2, 14.17, 17.13 y~20.1 (pandas baja el primero a 11.19 para incluir el mínimo), y los de la textura 14.1, 17.73, 21.37 y~25.0. Con la escala ordinal creciente para el radio y la textura, y la nominal para el diagnóstico, el contexto tiene $\lvert G \rvert = 6$ objetos y $\lvert M \rvert = 8$ atributos:

\begin{center}
\small
\begin{tabular}{c|ccc|ccc|cc}
  & \multicolumn{3}{c|}{radio} & \multicolumn{3}{c|}{textura} & \multicolumn{2}{c}{diagnóstico} \\
  & $\le 0$ & $\le 1$ & $\le 2$ & $\le 0$ & $\le 1$ & $\le 2$ & $=\mathrm{B}$ & $=\mathrm{M}$ \\ \hline
  $g_1$ & $\times$ & $\times$ & $\times$ & $\times$ & $\times$ & $\times$ & $\times$ &          \\
  $g_2$ & $\times$ & $\times$ & $\times$ &          & $\times$ & $\times$ & $\times$ &          \\
  $g_3$ &          &          & $\times$ &          & $\times$ & $\times$ &          & $\times$ \\
  $g_4$ & $\times$ & $\times$ & $\times$ &          &          & $\times$ & $\times$ &          \\
  $g_5$ &          &          & $\times$ &          & $\times$ & $\times$ &          & $\times$ \\
  $g_6$ &          & $\times$ & $\times$ & $\times$ & $\times$ & $\times$ & $\times$ &
\end{tabular}
\end{center}

Las columnas \enquote{radio${}\le 2$} y \enquote{textura${}\le 2$} están llenas (Corolario~\ref{cor:universal}); \enquote{radio${}\le 0$} está contenida en \enquote{radio${}\le 1$} (Proposición~\ref{prop:cadena}); las dos columnas del diagnóstico parten los objetos en $\{g_1, g_2, g_4, g_6\}$ y $\{g_3, g_5\}$ (Proposición~\ref{prop:particion}); y $g_3$ y~$g_5$ tienen la misma fila aunque sus valores originales son distintos, que es la pérdida de información de la discretización (Corolario~\ref{cor:filas}). El código que produce este ejemplo está en \path{docs/codigo/preprocessing/scaler.md}.

\section{Conceptos formales y NextClosure}
\label{sec:conceptos}

\emph{Pendiente: esta sección se completará en la Fase~2 (núcleo FCA).}


\section{Implicaciones y base de Guigues--Duquenne}
\label{sec:implicaciones}

\emph{Pendiente: esta sección se completará en la Fase~3 (implicaciones).}

\section{Dependencias funcionales y axiomas de Armstrong}
\label{sec:armstrong}

Las bases de datos relacionales usan, sobre la tabla original, una noción parecida a las implicaciones: las dependencias funcionales. Esta sección presenta los axiomas de Armstrong, que deducen implicaciones con reglas sintácticas y son equivalentes a la consecuencia de la \S\ref{sec:implicaciones-def}; el algoritmo de clausura que usa el código; la traducción entre dependencias funcionales e implicaciones, y el cálculo de claves candidatas. La clausura está implementada en \path{src/implications/armstrong.py}; las dependencias y las claves son la parte de la Fase~3 que falta implementar, y esta sección es su fundamento.


\subsection{Axiomas de Armstrong y clausura}
\label{sec:axiomas}

\begin{definicion}[reglas de Armstrong]\label{def:armstrong}
Para $X, Y, Z \subseteq M$: (R1)~\emph{reflexividad}: si $Y \subseteq X$, se obtiene $X \to Y$; (R2)~\emph{aumentación}: de $X \to Y$ se obtiene $X \cup Z \to Y \cup Z$; (R3)~\emph{transitividad}: de $X \to Y$ e $Y \to Z$ se obtiene $X \to Z$. Se escribe $\mathcal{L} \vdash X \to Y$ si hay una sucesión finita de implicaciones que termina en $X \to Y$ y en la que cada una está en~$\mathcal{L}$ o sale de las anteriores con una regla.
\end{definicion}

De ellas se derivan la unión (de $X \to Y$ y $X \to Z$, $X \to Y \cup Z$), la descomposición (de $X \to Y \cup Z$, $X \to Y$) y la pseudotransitividad (de $X \to Y$ y $W \cup Y \to Z$, $W \cup X \to Z$).

\begin{teorema}[\cite{armstrong1974}]\label{teo:armstrong}
$\mathcal{L} \vdash X \to Y$ si y solo si $\mathcal{L} \models X \to Y$, si y solo si $Y \subseteq \mathcal{L}(X)$.
\end{teorema}

\begin{proof}
La segunda equivalencia es la Proposición~\ref{prop:clausura-impl}. \emph{Corrección:} si $T$ respeta~$\mathcal{L}$, cada regla produce, a partir de implicaciones que $T$ respeta, otra que también respeta: (R1)~si $Y \subseteq X \subseteq T$, entonces $Y \subseteq T$; (R2)~si $X \cup Z \subseteq T$, entonces $Y \subseteq T$ y $Y \cup Z \subseteq T$; (R3)~si $X \subseteq T$, entonces $Y \subseteq T$ y $Z \subseteq T$. Por inducción, $T$ respeta todo lo que se deriva de~$\mathcal{L}$. \emph{Completitud:} con la sucesión~$W_s$ de la Proposición~\ref{prop:clausura-impl}, se ve que $\mathcal{L} \vdash X \to W_s$. Para $s = 0$ es (R1). Si $\mathcal{L} \vdash X \to U$, con $W_s \subseteq U$, y $(Z \to V) \in \mathcal{L}$ con $Z \subseteq W_s$, entonces $U \to Z$ por~(R1), $U \to V$ por~(R3), $U \to U \cup V$ por~(R2) y $X \to U \cup V$ por~(R3); haciendo esto con cada implicación aplicable se llega a $X \to W_{s+1}$. Así, $\mathcal{L} \vdash X \to \mathcal{L}(X)$, y si $Y \subseteq \mathcal{L}(X)$, (R1) y~(R3) dan $X \to Y$.
\end{proof}

La demostración es constructiva: cada cálculo de una clausura es una derivación con las reglas de Armstrong, en la que cada atributo agregado se justifica con una implicación de~$\mathcal{L}$. Por eso, en este proyecto, validar una implicación con los axiomas equivale a calcular una clausura.

\begin{teorema}[LinClosure, \cite{beeri1979}]\label{teo:linclosure}
Sea $W = X$ más las conclusiones de las implicaciones con premisa vacía. Cada implicación lleva un contador de los atributos de su premisa que faltan, cada atributo tiene la lista de implicaciones cuya premisa lo contiene, y una cola procesa una vez cada atributo de~$W$; cuando un contador llega a cero, los atributos nuevos de la conclusión se agregan a~$W$ y a la cola. El $W$ final es $\mathcal{L}(X)$, y se calcula con $O\bigl(\lvert M \rvert + \sum_{(Z \to V) \in \mathcal{L}} (\lvert Z \rvert + \lvert V \rvert)\bigr)$ operaciones.
\end{teorema}

\begin{proof}
Un atributo entra a la cola solo cuando se agrega a~$W$, así que entra una vez, y cada contador cuenta los atributos de su premisa que aún no han salido de la cola. $W \subseteq \mathcal{L}(X)$, porque una conclusión se agrega solo cuando su premisa ya está en~$W$, y $\mathcal{L}(X)$ respeta~$\mathcal{L}$. Al final la cola está vacía, así que toda implicación con premisa contenida en~$W$ llegó a contador cero, o tenía premisa vacía, y su conclusión está en~$W$: $W$ respeta~$\mathcal{L}$ y contiene a~$X$, luego $\mathcal{L}(X) \subseteq W$. Las listas suman $\sum \lvert Z \rvert$ y se recorren una vez, y cada conclusión se agrega a lo sumo una vez.
\end{proof}

\path{armstrong.py} tiene dos versiones. \texttt{attribute\_closure(X, F)} construye los contadores en cada llamada. \texttt{ImplicationClosure} recibe las implicaciones una por una con \texttt{add(A, B)}, se llama como una función, \texttt{ic(X)}, y solo crea contadores para las implicaciones que toca cada llamada; es la que usa Guigues--Duquenne, que calcula muchas clausuras mientras la lista crece. Las dos tratan aparte las premisas vacías, cuyo contador nunca bajaría, y las pruebas las comparan con la sucesión de punto fijo en 5000~casos aleatorios.


\subsection{Dependencias funcionales}
\label{sec:dependencias}

\begin{definicion}[dependencia funcional]\label{def:df}
Una \emph{relación}~$r$ sobre un conjunto finito~$R$ de columnas es un conjunto finito de tuplas $t \colon R \to W$, sin repeticiones ni valores faltantes. La \emph{dependencia funcional} $X \to Y$, con $X, Y \subseteq R$, \emph{se cumple} en~$r$ si dos tuplas que coinciden en~$X$ coinciden en~$Y$: $t[X] = u[X]$ implica $t[Y] = u[Y]$ \parencite{maier1983, abiteboul1995}.
\end{definicion}

Por la Observación~\ref{obs:relacional}, una relación es un contexto multivaluado; las dependencias hablan de la tabla original, y las implicaciones de la Sección~\ref{sec:implicaciones}, del contexto escalado. En un \texttt{DataFrame} puede haber filas repetidas, y hay que quitarlas antes (\texttt{drop\_duplicates}).

\begin{definicion}[contexto de pares]\label{def:pares}
El \emph{contexto de pares} de~$r$ es $\K_r = (P_r, R, I_{=})$, donde $P_r$ son los pares $\{t, u\}$ de tuplas distintas y $\{t, u\} \mathrel{I_{=}} a$ si $t[a] = u[a]$. La fila de un par es su \emph{conjunto de acuerdo} $\operatorname{ag}(t, u) = \{ a \in R : t[a] = u[a] \}$.
\end{definicion}

\begin{teorema}[\cite{ganter1999, baixeries2014}]\label{teo:df-implicaciones}
La dependencia funcional $X \to Y$ se cumple en~$r$ si y solo si la implicación $X \to Y$ es válida en~$\K_r$.
\end{teorema}

Las dos condiciones dicen lo mismo: todo par de tuplas cuyo conjunto de acuerdo contiene a~$X$ tiene un conjunto de acuerdo que contiene a~$Y$. En consecuencia, todo lo de las Secciones \ref{sec:conceptos} y~\ref{sec:implicaciones} se aplica a las dependencias: la clausura de~$X$ es $X^{+} = \bigcap \{ \operatorname{ag}(t, u) : X \subseteq \operatorname{ag}(t, u) \}$ (o $R$, si ningún par coincide en~$X$); la base de Guigues--Duquenne de~$\K_r$ es la lista más corta de dependencias de la que se deducen todas las de~$r$; y los axiomas de Armstrong son correctos y completos para las dependencias funcionales. Con la escala nominal, además, $X \to Y$ se cumple si y solo si, para cada tupla~$t$, la implicación $\{(a{=}t[a]) : a \in X\} \to \{(b{=}t[b]) : b \in Y\}$ es válida en el contexto derivado: una dependencia es la conjunción de una implicación entre valores por cada combinación de valores de~$X$.

En la práctica, $\K_r$ tiene $n(n-1)/2$ pares (161\,596 con $n = 569$), pero solo importan los conjuntos de acuerdo distintos, que se calculan agrupando las tuplas por el valor de cada columna \parencite{mannila1994}. Las dependencias dependen de la discretización: con intervalos más gruesos en~$X$, más pares coinciden en~$X$ y es más difícil que $X \to Y$ se cumpla. Y como las implicaciones, una dependencia se rompe con un solo par que no la cumpla.

\begin{ejemplo}[dependencias del ejemplo]\label{ej:df}
Se toma la tabla discretizada de la \S\ref{sec:ejemplo} como relación sobre $R = \{\mathrm{r}, \mathrm{t}, \mathrm{d}\}$ (radio, textura y diagnóstico). Las filas de $g_3$ y~$g_5$ son iguales, así que quedan cinco tuplas y diez pares, con conjuntos de acuerdo $\{\mathrm{r}, \mathrm{d}\}$ (tres pares: los de $g_1$, $g_2$ y~$g_4$), $\{\mathrm{t}, \mathrm{d}\}$ (el par $\{g_1, g_6\}$), $\{\mathrm{t}\}$ (el par $\{g_2, g_3\}$), $\{\mathrm{d}\}$ (dos pares) y~$\emptyset$ (tres pares). Todo par que coincide en el radio coincide en el diagnóstico, así que se cumple radio${}\to{}$diagnóstico, que es toda la base de~$\K_r$. En cambio, $\{\mathrm{t}\}^{+} = \{\mathrm{t}, \mathrm{d}\} \cap \{\mathrm{t}\} = \{\mathrm{t}\}$: la textura no determina nada.
\end{ejemplo}


\subsection{Claves candidatas}
\label{sec:claves}

\begin{definicion}[superclave y clave candidata]\label{def:clave}
Respecto de un conjunto~$\mathcal{F}$ de dependencias sobre~$R$, $X$ es una \emph{superclave} si $\mathcal{F}(X) = R$, y una \emph{clave candidata} si además es minimal; un atributo es \emph{primo} si está en alguna clave. Para una relación, $\mathcal{F}$ son las dependencias que se cumplen en ella y $\mathcal{F}(X) = X^{+}$.
\end{definicion}

\begin{proposicion}\label{prop:claves}
(1)~$X$ es una superclave de~$r$ si y solo si no hay dos tuplas distintas que coincidan en~$X$, es decir, si $X$ corta a todos los conjuntos $R \setminus \operatorname{ag}(t, u)$; las claves candidatas son los conjuntos minimales que los cortan a todos, y basta usar los conjuntos de acuerdo maximales. (2)~Si, empezando con $K = R$, se quita cada atributo~$a$ tal que $K \setminus \{a\}$ sigue siendo superclave, el resultado es una clave candidata, calculada con $\lvert R \rvert$ clausuras. (3)~Un atributo~$a$ está en todas las claves si y solo si $a \notin \mathcal{F}(R \setminus \{a\})$.
\end{proposicion}

La parte~(1) es la definición escrita con los conjuntos de acuerdo. Para~(2), un subconjunto de algo que no es superclave tampoco lo es, porque la clausura es monótona, así que ningún atributo que no se pudo quitar antes se puede quitar después. Para~(3), si $R \setminus \{a\}$ es superclave contiene una clave sin~$a$, y si no lo es, ningún conjunto sin~$a$ lo es. Una relación puede tener exponencialmente muchas claves: con $R = \{A_1, \dots, A_k, B_1, \dots, B_k\}$ y las dependencias $A_i \to B_i$ y $B_i \to A_i$ hay $2^k$. \textcite{lucchesi1978} dan un algoritmo para hallarlas todas en tiempo polinomial en el tamaño de la entrada y en el número de claves, y muestran que decidir si hay una clave de a lo sumo $q$~atributos es NP-completo.

\begin{ejemplo}[claves del ejemplo]\label{ej:claves}
En el Ejemplo~\ref{ej:df}, los conjuntos de acuerdo maximales son $\{\mathrm{r}, \mathrm{d}\}$ y $\{\mathrm{t}, \mathrm{d}\}$, con complementos $\{\mathrm{t}\}$ y~$\{\mathrm{r}\}$, así que la única clave candidata es $\{\mathrm{r}, \mathrm{t}\}$: radio y textura están en todas las claves, y el diagnóstico en ninguna. Sin quitar las filas repetidas, $g_3$ y~$g_5$ coincidirían en todo y la tabla no tendría ninguna clave; por eso se trabaja con la relación sin repeticiones, o se conserva una columna que identifique cada fila.
\end{ejemplo}

\section{Conceptos formales como características}
\label{sec:ml}

La fase opcional usa los conceptos del retículo como características binarias para clasificar con SVM y agrupar con $k$-medias, y compara los resultados con los de la representación original. Aprender a partir de conceptos tiene una larga tradición en FCA \parencite{kuznetsov2004ml}; aquí se usa la forma más simple, en la que cada concepto es una característica. El código de esta fase irá en \path{src/ml/pipeline.py}, con scikit-learn.

En esta sección, $G$ es el conjunto de objetos de entrenamiento, cada uno con su clase (en Breast Cancer Wisconsin, maligno o benigno). El contexto $\K = (G, M, I)$ se construye solo con las variables predictoras: la clase no es un atributo. Un objeto nuevo, discretizado y escalado con los mismos cortes y escalas, queda descrito por el conjunto $x \subseteq M$ de sus atributos binarios.


\subsection{Características de conceptos}
\label{sec:caracteristicas}

\begin{definicion}[características de conceptos]\label{def:caracteristicas}
Para una lista $\mathcal{C}$ de conceptos, la \emph{representación por conceptos} de un objeto~$x$ es el vector $\varphi(x) \in \{0, 1\}^{\mathcal{C}}$ con $\varphi_c(x) = 1$ si $B \subseteq x$ y $\varphi_c(x) = 0$ si no, para cada $c = (A, B) \in \mathcal{C}$. Para un objeto de entrenamiento, $x = g'$.
\end{definicion}

Con máscaras de bits, cada característica es una sola operación, $B \mathbin{\&} x = B$, y para un objeto nuevo no hay que recalcular el retículo.

\begin{proposicion}\label{prop:conjunciones}
(1)~Para $g \in G$, $\varphi_{(A, B)}(g') = 1$ si y solo si $g \in A$. (2)~Para todo $Y \subseteq M$, la función $g \mapsto [\, Y \subseteq g' \,]$ coincide en~$G$ con la característica del concepto $(Y', Y'')$, y conceptos distintos dan columnas distintas: las características de~$\Ret(\K)$ son exactamente las conjunciones de atributos, evaluadas en los objetos de entrenamiento, sin repeticiones. (3)~La característica del concepto atributo~$\mu m$ es la columna~$m$ de la matriz de incidencia.
\end{proposicion}

Las tres partes salen de que $B \subseteq g'$ equivale a $g \in B'$. Así, la representación por todos los conceptos contiene las columnas del contexto binario y les agrega sus conjunciones no redundantes; los experimentos miden si esas conjunciones ayudan. Para un objeto nuevo, $[\, Y \subseteq x \,]$ y $[\, Y'' \subseteq x \,]$ pueden diferir, porque la implicación $Y \to Y''$ de los datos de entrenamiento puede no cumplirse en~$x$; usar la intensión es la opción más específica, y usar un generador minimal, la más general.

Como el número de conceptos puede ser exponencial (con Animal Species hay 7\,387 para 200~objetos), hay que elegir una lista~$\mathcal{C}$ más corta. El \emph{soporte} de $(A, B)$ es $\sop(A, B) = \lvert A \rvert / \lvert G \rvert$, y el \emph{iceberg} de umbral~$\theta$ son los conceptos con soporte${}\ge \theta$ \parencite{stumme2002}. Si $c \le d$, entonces $\sop(c) \le \sop(d)$, así que el iceberg contiene a todos los conceptos por encima de los suyos; y como en Close-by-One la extensión de un hijo está contenida en la de su padre, la búsqueda se puede podar en cuanto el soporte baja de~$\theta$. Otra opción es la \emph{estabilidad} de un concepto, la fracción de subconjuntos $C \subseteq A$ con $C' = B$ \parencite{kuznetsov2007}, que favorece a los conceptos que no dependen de unos pocos objetos.

\begin{proposicion}[núcleo del retículo]\label{prop:nucleo}
Con $\mathcal{C} = \Ret(\K)$ y $g, h \in G$, $\langle \varphi(g'), \varphi(h') \rangle$ es el número de conceptos por encima de $\gamma g \vee \gamma h$, que es el número de intensiones contenidas en $g' \cap h'$; y $\lVert \varphi(g') - \varphi(h') \rVert^2$ es el número de conceptos cuya extensión contiene a uno solo de los dos.
\end{proposicion}

Se sigue de las Proposiciones \ref{prop:objeto-atributo} y~\ref{prop:conjunciones}: un concepto contiene a $g$ y a~$h$ si y solo si está por encima de $\gamma g$ y de~$\gamma h$. Dos objetos son parecidos si comparten muchos conceptos, y una SVM lineal sobre estas características es una SVM con este núcleo.


\subsection{Modelos}
\label{sec:modelos}

Sean $(x_i, y_i)$, con $x_i \in \R^d$ e $y_i \in \{-1, +1\}$, los datos de entrenamiento. La \emph{máquina de vectores de soporte} de margen suave \parencite{cortes1995} busca el clasificador $f(x) = \operatorname{signo}(\langle w, x \rangle + b)$ que resuelve
\[
  \min_{w, b} \ \frac{1}{2} \lVert w \rVert^2 + C \sum_{i} \max\bigl( 0,\ 1 - y_i (\langle w, x_i \rangle + b) \bigr).
\]
El primer término maximiza el margen, que es $2/\lVert w \rVert$, y el segundo penaliza, con peso~$C$, los datos que quedan dentro del margen o del lado equivocado (la pérdida bisagra). La solución es $w = \sum_i \alpha_i y_i x_i$, donde solo los \emph{vectores de soporte} tienen $\alpha_i > 0$, y depende de los datos solo a través de los productos internos $\langle x_i, x_j \rangle$, que se pueden reemplazar por un núcleo \parencite{cortes1995, hastie2009}. Sobre las características de conceptos, cada peso~$w_c$ corresponde a una conjunción de condiciones sobre intervalos, como \enquote{radio${}\le 17.13$ y textura${}\le 21.37$}, así que los pesos más grandes se leen como reglas. Las variables numéricas originales se estandarizan con la media y la desviación de la parte de entrenamiento, y $C$ se elige por validación cruzada.

El algoritmo de Lloyd en~$\R^d$ tiene la misma garantía que en una dimensión: el Teorema~\ref{teo:lloyd} vale palabra por palabra, aplicando el Lema~\ref{lem:media} a cada coordenada. scikit-learn elige los centroides iniciales con $k$-means++ \parencite{arthur2007} y repite el algoritmo con varias inicializaciones, porque el resultado es un óptimo local (Observación~\ref{obs:optimo-local}). Para vectores binarios, $\lVert u - v \rVert^2$ es el número de coordenadas en que difieren, y cada coordenada de un centroide es la fracción de objetos del grupo con esa característica.


\subsection{Evaluación sin fuga de información}
\label{sec:evaluacion}

Con la clase maligna como positiva y VP, FP, VN y FN los verdaderos y falsos positivos y negativos,
\begin{gather*}
  \text{exactitud} = \frac{\mathrm{VP} + \mathrm{VN}}{n}, \qquad
  \text{sensibilidad} = \frac{\mathrm{VP}}{\mathrm{VP} + \mathrm{FN}}, \\
  \text{especificidad} = \frac{\mathrm{VN}}{\mathrm{VN} + \mathrm{FP}}, \qquad
  F_1 = \frac{2\,\mathrm{VP}}{2\,\mathrm{VP} + \mathrm{FP} + \mathrm{FN}}.
\end{gather*}
En diagnóstico, el error más costoso es el falso negativo, así que la sensibilidad se reporta junto con la exactitud. Los modelos se evalúan con \emph{validación cruzada estratificada}: los objetos se dividen en $F$~partes con la misma proporción de clases, cada parte se predice con un modelo entrenado en las demás, y se reportan la media y la desviación de cada métrica \parencite{hastie2009}.

\begin{observacion}[qué se ajusta dentro de cada partición]\label{obs:fuga}
Todo lo que mira los datos se calcula solo con la parte de entrenamiento y se aplica, sin cambios, a la de prueba: el número de intervalos y los cortes de la discretización (con los cortes extremos en $-\infty$ y $+\infty$, para que los valores nuevos fuera de rango caigan en el primer o el último intervalo), los valores de cada escala, el retículo y los conceptos elegidos, la estandarización y los hiperparámetros, que se eligen con una segunda validación cruzada dentro del entrenamiento. Calcular el retículo con todos los datos antes de dividirlos sería una fuga, y si la clase fuera un atributo del contexto, sus conceptos atributo serían características iguales a la clase.
\end{observacion}

Para el agrupamiento se usan el \emph{índice de Rand ajustado} \parencite{hubert1985}, que compara los grupos con las clases y vale~1 si coinciden y~0 en promedio si son aleatorios, y la \emph{silueta media} \parencite{rousseeuw1987}, que mide qué tan separados están los grupos sin usar las clases: para cada objeto, $s = (b - a)/\max(a, b)$, con $a$ su distancia media a su grupo y $b$ la menor distancia media a otro grupo.

\begin{observacion}[diseño del experimento]\label{obs:diseno}
Se comparan tres representaciones: (R1)~las variables originales estandarizadas, (R2)~las columnas del contexto binario y (R3)~las características de conceptos, de todo el retículo o de un iceberg. Para clasificar se usa una SVM lineal en las tres, y una SVM con núcleo gaussiano en~R1 como referencia, con validación cruzada estratificada de 10~partes; se reportan exactitud, sensibilidad, especificidad, $F_1$, número de características y tiempo, comparando partición por partición. Para agrupar se usa $k$-medias con $k = 2$ y se reportan el índice de Rand ajustado respecto del diagnóstico y la silueta. Como R3 contiene a~R2 (Proposición~\ref{prop:conjunciones}), la comparación entre las dos mide lo que aportan las conjunciones.
\end{observacion}

\begin{ejemplo}[características del ejemplo]\label{ej:caracteristicas}
Sin las columnas del diagnóstico, el contexto de los tumores tiene 8~conceptos, con intensiones $c_0 = \{r_2, t_2\}$, $c_1 = \{r_2, t_1, t_2\}$, $c_2 = \{r_1, r_2, t_2\}$, $c_3 = \{r_1, r_2, t_1, t_2\}$, $c_4 = \{r_1, r_2, t_0, t_1, t_2\}$, $c_5 = \{r_0, r_1, r_2, t_2\}$, $c_6 = \{r_0, r_1, r_2, t_1, t_2\}$ y $c_7 = M$. La representación por conceptos es:
\begin{center}
\small
\begin{tabular}{c|cccccccc|c}
  & $c_0$ & $c_1$ & $c_2$ & $c_3$ & $c_4$ & $c_5$ & $c_6$ & $c_7$ & Clase \\ \hline
  $g_1$ & 1 & 1 & 1 & 1 & 1 & 1 & 1 & 1 & B \\
  $g_2$ & 1 & 1 & 1 & 1 & 0 & 1 & 1 & 0 & B \\
  $g_3$ & 1 & 1 & 0 & 0 & 0 & 0 & 0 & 0 & M \\
  $g_4$ & 1 & 0 & 1 & 0 & 0 & 1 & 0 & 0 & B \\
  $g_5$ & 1 & 1 & 0 & 0 & 0 & 0 & 0 & 0 & M \\
  $g_6$ & 1 & 1 & 1 & 1 & 1 & 0 & 0 & 0 & B
\end{tabular}
\end{center}
La columna de~$c_2$, el concepto atributo de~$r_1$, separa las dos clases: \enquote{radio${}\le 17.13$} identifica a los benignos de esta muestra. Por la Proposición~\ref{prop:nucleo}, $g_1$ y~$g_2$ comparten seis conceptos y $g_1$ y~$g_3$ solo dos. $g_3$ y~$g_5$ tienen el mismo vector, porque la discretización los volvió indistinguibles (Corolario~\ref{cor:filas}).
\end{ejemplo}


\printbibliography[heading=bibintoc, title={Referencias}]

\end{document}
